\section{L'impasse ergonomique}
Si l'infrastructure technique de CujasNum pose d'importants défis en termes de sécurité, la limite se révèle dans son fonctionnement quotidien. La chaîne de numérisation actuelle de la bibliothèque numérique illustre parfaitement l'approche décrite par Elina Leblanc dans sa thèse la qualifiant d'une \enquote{approche purement descendante(top-down), axée sur les contraintes du système plutôt que sur l'activité des agents} \footcite{leblancBibliothequesNumeriquesEnrichies}.

Elina Leblanc s'appuie sur les critères d'Aaron Schmidt, l'évaluation d'un service de bibliothèque numérique présente \enquote{un guide pour la CCU des bibliothèques analogiques [...]sous un angle utilisateur} \footcite{leblancBibliothequesNumeriquesEnrichies}. Il mentionne ainsi l'utilité d'un service d'une bibliothèque numérique en disant que celui-ci doit être \enquote{utilisable et désirable}. Aaron Schmidt l'explique ainsi : \enquote{Un service utile doit répondre à un besoin et avoir une importance aux yeux des utilisateurs pour qu'ils ne l'utilisent pas de manière ponctuelle, mais sur la longue durée. Un service utilisable garantit aux utilisateurs un accès aisé et rapide aux données. Enfin, la désirabilité crée un lien entre l'utilisateur et un service, en donnant envie au premier d'utiliser le deuxième et en suscitant un besoin} \footcite{leblancBibliothequesNumeriquesEnrichies}. Cela garantit à l'utilisateur une expérience optimale lorsqu'il utilise la bibliothèque numérique d'une institution patrimoniale. Or, à la Bibliothèque Cujas, l'interface n'est pas adaptée à la manière dont l'humain utilise les sites web aujourd'hui. L'utilisateur doit s'adapter aux contraintes rigides de l'interface et à la façon dont lui sont présentés les services. 

Ce manque d'ergonomie se traduit par un décalage dans la production des ouvrages numériques. Si le contrôle qualité d'un ouvrage requiert environ \enquote{quinze minutes en moyenne si tout va bien}, la moindre anomalie structurelle ou typographique fait basculer le traitement dans une autre dimension nécessitant au minimum \enquote{trois jours} de manipulations afin de corriger l'erreur détectée et ensuite pouvoir remettre en ligne le document \footannexe{\nameref{annexe_A_rapport_audit}.}. Cette manière de travailler est devenue un gouffre chronophage et la rigidité qu'implique ce traitement est devenue un carcan dont la Bibliothèque Cujas ne parvient pas à se défaire.

L'architecture d'\ac{XTF}, pensée pour les monographies, ne permet pas de l'adapter à d'autres types de documents avec des arborescences complexes, comme les périodiques. \ac{XTF} repose sur une arborescence dite \enquote{à plat}\footannexe{\nameref{annexe_A_rapport_audit}.}. C'est la raison pour laquelle si on ajoute d'autres types de documents l'arobrescence ne sera pas respectée, constituant un sérieux problème pour les équipes de Cujas qui souhaitent ajouter d'autres collections dans CujasNum.

Contrairement à une monographie, un périodique exige une architecture d'accès d'au moins trois niveaux, \enquote{c'est-à-dire le titre de la revue, suivie de l'année de publication et du numéro du volume ou du fascicule, pour que la consultation soit pertinente} \footannexe{\nameref{annexe_A_rapport_audit}.}. Etant incapable de reproduire cette structure, c'est l'une des raisons qui pousse la bibliothèque Cujas à mettre de côté ses collections ne pouvant les intégrer sans \enquote{dénaturer} le type de document. Aujourd'hui, on observe sur certaines bibliothèques numériques possédant des périodiques que la granularité peut aller jusqu'à l'article même. \ac{XTF} s'arrête initialement au volume et ne permet pas une granularité aussi précise, allant donc à l'encontre des pratiques de recherche actuelle. Faute de pouvoir diffuser ces revues dans sa bibliothèque numérique, Cujas a fait \enquote{le choix de faire numériser ces volumes via un partenariat avec Persée} \footannexe{\nameref{annexe_A_rapport_audit}.}. Pour signaler ces numérisations, la bibliothèque se contente uniquement de créer des notices de redirection indiquant aux chercheurs l'existance de ces documents, où il peut les trouver afin de les consulter. 

Un autre problème réside dans l'indexation utilisée par la bibliothèque numérique actuelle, l'indexation propre à la bibliothèque Cujas. Celle-ci ne permet pas le filtrage via différents types de documents; en effet elle a été structurée autour de la monographie ne permettant pas de croiser les informations \footcite{2013_Num_Indexation_manuel_FHdocEspaceTravail}. L'impossibilité d'ajouter d'autres types de documents met en lumière l'invisibilisation à laquelle on se trouve confronté lors d'une recherche plein-texte. Ces documents ne peuvent ressortir de manière évidente ne pouvant posséder d'\ac{OCR}. Les résultats apparaissent donc de façon très restreinte et dans les dernières pages problablement \footannexe{\nameref{annexe_A_rapport_audit}.}. L'obsolescence du logiciel \ac{XTF} entraîne donc \textit{de facto} un filtrage empêchant la Bibliothèque Cujas de pouvoir mettre en avant la richesse et la diversité de ses collections. Il lui est impossible de valoriser ses documents afin de justifier les choix effectués ou de mettre en avant une collection ou un document dans un article ou dans une exposition virtuelle. 
\medskip

Le bilan de CujasNum met en avant les points forts du projet ainsi que les points faibles au moment de la réalisation de celui-ci, mais également ceux ayant émergé par la suite avec l'évolution des technoloqies numériques et des pratiques du Web. D'abord, pionnière dans le monde des bibliothèques numériques, la Bibliothèque Cujas s'est trouvée progressivement isolée, parvenant difficilement à suivre l'évolution rapide liée à l'émergence de l'internet entraînant les pratiques documentaires et les pratiques d'Internet dans un nouvel âge. Face à ce constat, \enquote{l'ajout d'une surcouche WordPress ne s'aurait être la solution comme cela avait été envisagé dans un premier temps} \footannexe{\nameref{annexe_A_rapport_audit}.}. Pour répondre aux attentes contemporaines des chercheurs et pour s'adapter aux mutations nationales, des systèmes de gestion des bibliothèques (SGB) imposent une rupture dans les pratiques actuelles. Cela nécessite de réfléchir à une interface plus souple, acceptant différents types de documents se basant sur les pratiques actuelles de la recherche, répondant aux nouveaux enjeux des utilisateurs. La bibliothèque numérique doit permettre à sa bibliothèque d'utiliser les standards du Web de données afin de rendre sa plateforme davantage interopérable et prête à s'intégrer dans de grands réseaux scientifiques nationaux.